% Duan HOFAS / sub-FAS IEEEtran LaTeX template.
% Replace every placeholder and verify all mathematics against the selected source.
\IfFileExists{IEEEtran.cls}
  {\documentclass[journal]{IEEEtran}}
  {\documentclass[11pt]{article}}

\usepackage{amsmath,amssymb,amsfonts,amsthm,bm,booktabs,cite,graphicx}

\newtheorem{definition}{Definition}
\newtheorem{assumption}{Assumption}
\newtheorem{problem}{Problem}
\newtheorem{lemma}{Lemma}
\newtheorem{proposition}{Proposition}
\newtheorem{theorem}{Theorem}
\newtheorem{corollary}{Corollary}
\theoremstyle{remark}
\newtheorem{remark}{Remark}
\newtheorem{example}{Example}

\newcommand{\R}{\mathbb{R}}
\newcommand{\xstack}{x^{(0\sim n-1)}}
\newcommand{\Arow}{A_{0\sim n-1}}
\newcommand{\Phic}{\Phi(\Arow)}

\begin{document}

\title{Replace With a Precise HOFAS or Sub-FAS Title}
\author{Author Names}
\maketitle

\begin{abstract}
% 5-part Duan abstract: background/limitation -> system class & core concept ->
% controller/switching construction -> closed-loop & ROEA property -> examples.
State the problem limitation, system class, proposed HOFAS/sub-FAS method,
principal closed-loop/ROEA results, and illustrative examples.
\end{abstract}

\section{Introduction}
% Target ~1.5 pages for an ~8-page IEEE Transactions manuscript.
% Flow: background & state-space bottleneck -> HOFAS/sub-FAS literature & gap ->
% core idea -> contribution list -> organization -> italic Notations paragraph.

The main results of this paper are summarized as follows.
\begin{enumerate}
  \item State the first structural/modeling contribution.
  \item State the second controller/switching contribution.
  \item State the third closed-loop/ROEA characterization contribution.
\end{enumerate}

The remainder of this paper is organized as follows. Section~II introduces the
system model and preliminaries. Section~III presents the controller design and
main results. Section~IV gives numerical examples, and Section~V concludes the
paper.

\emph{Notations:} For simplicity, the argument \(t\) of a time-dependent
variable is omitted when no confusion arises. For positive integers \(p\) and
\(q\), \(\R^p\) and \(\R^{p\times q}\) denote the \(p\)-dimensional real vector
space and the set of all \(p\times q\) real matrices, respectively. The symbols
\(I_p\) and \(0_{p\times q}\) denote the identity and zero matrices of the
indicated dimensions. For \(x\in\R^r\), \(x^{(j)}\) denotes its \(j\)th
derivative, and
\begin{equation*}
  \xstack
  =\begin{bmatrix}
     x^{\mathrm{T}} & \dot x^{\mathrm{T}} & \cdots & \bigl(x^{(n-1)}\bigr)^{\mathrm{T}}
   \end{bmatrix}^{\mathrm{T}}\in\R^{nr}.
\end{equation*}
For matrices \(A_j\in\R^{r\times r}\), \(j=0,1,\ldots,n-1\), define
\begin{equation*}
  \Arow=[\,A_0\ A_1\ \cdots\ A_{n-1}\,]
\end{equation*}
and
\begin{equation*}
  \Phic
  =\begin{bmatrix}
     0_{r\times r}&I_r&\cdots&0_{r\times r}\\
     0_{r\times r}&0_{r\times r}&\ddots&\vdots\\
     \vdots&\vdots&\ddots&I_r\\
     -A_0&-A_1&\cdots&-A_{n-1}
   \end{bmatrix}.
\end{equation*}

\section{Model Description and Preliminaries}

Consider a single-order affine system in the following form:
\begin{equation}\label{eq:model}
  x^{(n)}=f\bigl(\xstack,t\bigr)+B\bigl(\xstack,t\bigr)u,
\end{equation}
where \(n\geq 1\) is an integer, \(x\in\R^r\) and \(u\in\R^r\) are the state
vector and the control input vector, respectively, and \(f(\cdot,t)\in\R^r\)
and \(B(\cdot,t)\in\R^{r\times r}\) are piecewise continuous vector and matrix
functions, respectively.

\begin{definition}\label{def:feasible-set}
State the source-faithful singular set \(\mathcal{S}\), feasible set
\(\mathcal{F}\), or class definition here.
\end{definition}

\begin{assumption}\label{ass:main}
State only the exact additional assumption needed by the analysis, with a clear
use-site in the proof.
\end{assumption}

\section{Controller Design and Main Results}

When \(\mathcal{F}\neq\emptyset\), \(B^{-1}\) exists on \(\mathcal{F}\). Hence,
we construct the following controller on \(\mathcal{F}\):
\begin{equation}\label{eq:controller}
  u=-B^{-1}\bigl(\xstack,t\bigr)\bigl[\Arow\xstack+f\bigl(\xstack,t\bigr)-v\bigr],
\end{equation}
where \(A_j\in\R^{r\times r}\), \(j=0,1,\ldots,n-1\), are design matrices.

% Every Lemma, Proposition, Theorem, and Corollary MUST state its purpose
% explicitly inside its own statement.
\begin{theorem}\label{thm:main}
Suppose that system~\eqref{eq:model} satisfies Assumption~\ref{ass:main}. Let
the free parameters be arbitrarily given, and let \(A_j\in\R^{r\times r}\),
\(j=0,1,\ldots,n-1\), be a set of matrices such that \(\Phic\) is Hurwitz.
Consider the control law~\eqref{eq:controller} with \(v=0\). Then state the
explicit closed-loop stability, convergence, and ROEA characterization here.
\end{theorem}

\begin{proof}
Substituting control law~\eqref{eq:controller} into system~\eqref{eq:model}
gives the closed-loop system
\begin{equation}\label{eq:closed-loop}
  x^{(n)}+\Arow\xstack=0,
\end{equation}
which is equivalent to the following state-space form:
\begin{equation}\label{eq:state-space}
  \dot{x}^{(0\sim n-1)}=\Phic\xstack.
\end{equation}
% Complete the Lyapunov / trajectory-feasibility / contradiction argument.
% Contradiction template (Duan & Zhou 2022 Lemma 4.4):
%   Suppose that ... Then it follows from (...) that ... Therefore, there
%   exists ... such that ..., which means that ... On the other side, ...
%   This yields a contradiction. Thus ... The proof is finished.
The proof is finished.
\end{proof}

\begin{remark}
Discuss design degrees of freedom, special cases, or implementation bounds here
without hiding theorem hypotheses or proof-bearing main claims.
\end{remark}

\section{Numerical Examples}

\begin{table}[!t]
  \caption{System Parameters}
  \label{tab:params}
  \centering
  \begin{tabular}{ccc}
    \toprule
    Symbol & Value & Meaning \\
    \midrule
    \(a_0\) & \(4\) & stiffness coefficient \\
    \(a_1\) & \(1.2\) & damping coefficient \\
    \bottomrule
  \end{tabular}
\end{table}

% Use Duan narration: "The simulation results are shown in Fig. 1. It can be
% observed that ..." and short 2-6 word noun-phrase captions.

\section{Conclusion}
Summarize the problem addressed, the main closed-loop and ROEA results, the
numerical verification, and unaddressed scope boundaries.

\end{document}
